\SourceLeaf{423}%
durch die $3$ Punkte $(1),(2),(3)$ bezeichnet werden, die also je $90^{0}$ von einander entfernt sind, und zugleich den Sinn bezeichnen sollen, in welchem die Coordinaten als wachsend betrachtet werden. Wir stellen hier einige bekannte Sätze, von denen beständig Gebrauch gemacht werden wird, zusammen.

1) Der Winkel zweier sich schneidenden geraden Linien wird durch den Bogen [des grössten Kreises] zwischen den Punkten gemessen, die auf der Kugelfläche ihre Richtung vorstellen.

2) Die Lage einer jeden Ebene kann auf der Kugelfläche durch den grössten Kreis vorgestellt werden, in welchem jene von einer mit ersterer Ebene parallel durch den Mittelpunkt der Kugel gelegten Ebene geschnitten wird.

3) Der Winkel zweier Ebenen ist gleich dem Winkel der grössten Kreise, die die Lage von jenen repräsentiren, und wird also auch durch den Winkel zwischen den Polen der grössten Kreise gemessen.

4) Sind $x,y,z$; $x',y',z'$ die Coordinaten zweier Punkte, $r$ ihre Entfernung von einander und $L$ der Punkt der Kugelfläche, der die Richtung der geraden Linie vom ersten zum zweiten vorstellt, so ist
\[
\begin{aligned}
x'&=x+r\cos(1)L\\
y'&=y+r\cos(2)L\\
z'&=z+r\cos(3)L.
\end{aligned}
\]
\par\noindent\hspace*{1.2em}5) Es folgt hieraus leicht, dass allemal
\[
\cos(1)L^{2}+\cos(2)L^{2}+\cos(3)L^{2}=1
\]
[und] auch, wenn $L'$ irgend ein anderer Punkt der Kugelfläche ist,
\[
\cos(1)L.\cos(1)L'+\cos(2)L.\cos(2)L'+\cos(3)L.\cos(3)L'=\cos LL'.
\]

Wir fügen noch ein anderes Theorem bei, welches unseres Wissens sonst nirgends vorkommt und öfters mit Nutzen gebraucht werden kann.

Es seien $L,L',L'',L'''$ vier Punkte der Kugelfläche, und $A$ der Winkel, welchen $LL'''$, $L'L''$ an ihrem Durchschnittspunkte bilden. [Dann hat man]
\[
\cos LL'.\cos L''L'''-\cos LL''.\cos L'L'''=\sin LL'''.\sin L'L''.\cos A.
\]
